\documentclass[aspectratio=169, 14pt]{beamer}
\usepackage{stampfly_slides}

\lesson{12}{精密着陸競技会}{Precision Landing Competition}

\begin{document}

% ------------------------------------------------------------------
\begin{frame}
  \titlepage
\end{frame}

% ------------------------------------------------------------------
\begin{frame}{競技会概要 / Competition Overview}
  \begin{block}{精密着陸競技}
    パイロットは定位置から操縦し、3\,m 先のヘリポートに精密着陸する
  \end{block}

  \vspace{1em}

  \begin{itemize}
    \item \textbf{種目}: 精密着陸（パイロット定位置、3\,m 先のヘリポートに着陸）
    \item \textbf{ヘリポート}: 約 40\,cm $\times$ 40\,cm
    \item \textbf{評価}: 着陸までのタイム（ベストタイム採用）
  \end{itemize}
\end{frame}

% ------------------------------------------------------------------
\begin{frame}{競技ルール詳細 / Competition Rules}
  \centering
  \begin{tabular}{ll}
    \hline
    \textbf{項目} & \textbf{内容} \\
    \hline
    種目      & 精密着陸（パイロット定位置操縦） \\
    距離      & パイロットからヘリポートまで 3\,m \\
    ヘリポート & 40\,cm $\times$ 40\,cm \\
    制限時間  & 60 秒 \\
    試行回数  & 3 回（ベストタイム採用） \\
    判定      & ARM → 離陸 → ヘリポート着陸までの時間 \\
    パイロット & 定位置から動かない \\
    失格条件  & 安全エリア外飛行、手による介入、パイロットの移動 \\
    \hline
  \end{tabular}

  \vspace{1em}

  \begin{alertblock}{注意}
    安全第一！\ 周囲に人がいないことを確認してから飛行
  \end{alertblock}
\end{frame}

% ------------------------------------------------------------------
\begin{frame}{タイムスケジュール / Schedule}
  \centering
  \begin{tabular}{ll}
    \hline
    \textbf{時間} & \textbf{内容} \\
    \hline
    \multicolumn{2}{l}{\textbf{Day 4（本日午後）}} \\
    13:00--13:30 & ルール説明（このスライド） \\
    13:30--16:00 & 準備・自由練習・予選 \\
    \hline
    \multicolumn{2}{l}{\textbf{Day 5（最終日）}} \\
    \phantom{0}9:00--\phantom{0}9:15 & ルール確認・機体チェック \\
    \phantom{0}9:15--\phantom{0}9:45 & 最終チューニング \\
    \phantom{0}9:45--10:45 & 競技本番 \\
    10:45--11:30 & 結果発表・表彰・振り返り \\
    \hline
  \end{tabular}
\end{frame}

% ------------------------------------------------------------------
\begin{frame}{攻略のヒント / Tips for Success}
  \begin{columns}[t]
    \begin{column}{0.48\textwidth}
      \begin{exampleblock}{PID ゲイン調整}
        \begin{itemize}
          \item 安定性重視（急旋回よりも穏やかな応答）
          \item ヨー制御で方向を合わせる
          \item スロットル微調整で高度を安定させる
        \end{itemize}
      \end{exampleblock}
    \end{column}
    \begin{column}{0.48\textwidth}
      \begin{exampleblock}{練習のポイント}
        \begin{itemize}
          \item まず安定ホバリングを確立
          \item 前進 → 停止 → 降下の手順を練習
          \item バッテリ・プロペラを毎回チェック
        \end{itemize}
      \end{exampleblock}
    \end{column}
  \end{columns}

  \smallskip
  {\small \textbf{テレメトリ活用:}
    Teleplot でリアルタイムデータを確認し、ゲインを追い込む}
\end{frame}

% ------------------------------------------------------------------
\begin{frame}{チェックポイント / Checkpoint}
  \begin{alertblock}{準備確認}
    \begin{itemize}
      \item[$\square$] 競技ルールと失格条件を理解した
      \item[$\square$] バッテリがフル充電されている
      \item[$\square$] PID ゲインの調整方法を把握している
      \item[$\square$] 安全チェックリストを確認した
    \end{itemize}
  \end{alertblock}

  \vspace{1em}

  \begin{block}{さあ、練習を始めよう！}
    13:30 から自由練習 → 最高の精密着陸を目指そう
  \end{block}
\end{frame}

\end{document}
